% ----------------------------------------------------------------------
%  3.1.5. Типы связи: обзор и ковалентная связь
% ----------------------------------------------------------------------
\subsection{Типы связи в твёрдых телах}\label{ssec:bond-types}

\subsubsection{Четыре типа твёрдых тел и четыре типа связи}

Строгого определения типов твёрдых тел нет, но по совокупности свойств
кристаллы обычно делят на четыре группы (табл.~\ref{tab:solid-types}).
Каждой группе отвечает свой преобладающий тип связи.

\begin{table}[ht]
\centering
\caption{Типы кристаллических твёрдых тел и связи}
\label{tab:solid-types}
\small
\begin{tabularx}{\textwidth}{@{}>{\raggedright\arraybackslash}p{24mm}>{\raggedright\arraybackslash}p{30mm}>{\raggedright\arraybackslash}X>{\raggedright\arraybackslash}p{24mm}@{}}
\toprule
Тип кристалла & Из каких атомов & Характерные свойства & Связь \\
\midrule
Металлы & электро\-положи\-тельные элементы & пластичны, высокая электро- и теплопроводность (носители~— электроны), непрозрачны, металлический блеск; $\Tpl$ в широких пределах & металли\-ческая («$+$» с «$+$») \\
Ионные & сильно электро\-положи\-тельные + сильно электро\-отрица\-тельные & ионная проводимость, растущая с температурой; свободных электронов нет~— прозрачны & ионная («$+$» с «$-$») \\
Валентные (кова\-лентные) & лёгкие элементы средних групп & нет ионной и (обычно) электронной проводимости, при нагреве~— полупроводники; твёрдые и хрупкие & ковалентная («$-$» с «$-$») \\
Молеку\-лярные & благородные газы, насыщенные молекулы & низкие $\Tpl$ и $T_\text{кип}$, испаряются в виде молекул, не проводят ток & ван-дер-ваальсова («0» с «0») \\
\bottomrule
\end{tabularx}
\end{table}

Деление условно: есть вещества, которые нельзя отнести только к одному
типу. Пример~— \term{графит} (рис.~\ref{fig:graphite}). Он состоит из
слоёв, в каждом из которых атом углерода связан с тремя соседями на
расстоянии $0{,}141$~нм прочными ковалентными связями: каждый слой~—
огромная плоская молекула. Слои отстоят друг от друга на $0{,}34$~нм и
связаны слабо, как молекулы в молекулярном кристалле; четвёртый
валентный электрон каждого атома делокализован и даёт связь, похожую на
металлическую. Отсюда сочетание свойств: прочность вдоль слоёв, лёгкое
расслаивание, металлический блеск и электронная проводимость.

\begin{figure}[ht]
\centering
\begin{tikzpicture}[x={(1cm,0cm)}, y={(0.5cm,0.28cm)}, z={(0cm,1cm)}, scale=0.62]
  \newcommand{\hexlayer}[3]{% z, сдвиг по x, цвет
    \foreach \i in {0,...,3}{\foreach \j in {0,...,3}{
      \pgfmathsetmacro{\cx}{\i*1.732+\j*0.866+#2}
      \pgfmathsetmacro{\cy}{\j*1.5}
      \foreach \a in {30,90,...,330}{
        \draw[#3, line width=0.6pt] ({\cx+cos(\a)},{\cy+sin(\a)},#1) -- ({\cx+cos(\a+60)},{\cy+sin(\a+60)},#1);}}}}
  \hexlayer{0}{0}{atomLine!70}
  \hexlayer{2.3}{0.866}{fmBlue}
  \draw[dimline] (8.6,0,0) -- node[right, slbl]{$0{,}34$~нм} (8.6,0,2.3);
\end{tikzpicture}
\caption{Слоистая структура графита (схема): в каждом слое атомы образуют
гексагональную сетку (расстояние между соседями $0{,}141$~нм), соседние
слои сдвинуты и отстоят друг от друга на $0{,}34$~нм}
\label{fig:graphite}
\end{figure}

\subsubsection{Ковалентная связь}

\paragraph{Ион молекулы водорода H$_2^+$}
Начнём с простейшей системы~— два протона и один электрон (такие ионы
действительно существуют, их наблюдают масс-спектрометрически). Пока атом
водорода и протон далеко (5~нм), взаимодействия нет. При сближении до
$\sim0{,}2$~нм электрон начинает «перескакивать» от одного протона к
другому, а при равновесном расстоянии $0{,}106$~нм он становится
\emph{общим}. Главное: электронное облако втягивается в пространство
\emph{между} протонами и стягивает их, причём ни одно ядро не имеет
преимущества перед другим.

\paragraph{Молекула H$_2$}
Так же образуется молекула водорода, только общими становятся два
электрона~— по одному от каждого атома, с противоположными спинами
(рис.~\ref{fig:H2}). Ковалентная связь~— \emph{парная}: её образует пара
электронов, концентрирующих высокую плотность заряда между ядрами
(длина связи H--H $0{,}074$~нм, энергия $\approx4{,}5\eV$).

\begin{figure}[ht]
\centering
\begin{tikzpicture}[scale=0.85]
  \begin{scope}
    \foreach \r/\o in {1.2/0.10, 0.95/0.12, 0.7/0.15, 0.45/0.2}{
      \fill[fmBlue, opacity=\o] (-1.5,0) circle (\r);}
    \fill[fmRed] (-1.5,0) circle (2.2pt);
    \fill[fmRed] (1.5,0) circle (2.2pt);
    \node[slbl, below=3pt] at (1.5,0) {$p^+$};
    \node[slbl] at (0,-1.6) {атом H и протон};
  \end{scope}
  \begin{scope}[shift={(6,0)}]
    \foreach \sx/\sy/\o in {1.45/0.95/0.12, 1.2/0.78/0.14, 0.95/0.6/0.17, 0.7/0.45/0.22}{
      \fill[fmBlue, opacity=\o] (0,0) ellipse ({\sx} and {\sy});}
    \foreach \x in {-0.53,0.53}{\fill[fmBlue, opacity=0.12] (\x,0) circle (0.5);}
    \fill[fmRed] (-0.53,0) circle (2.2pt);
    \fill[fmRed] (0.53,0) circle (2.2pt);
    \node[slbl] at (0,-1.6) {молекула H$_2$};
  \end{scope}
\end{tikzpicture}
\caption{Образование ковалентной связи: электронное облако (синий цвет)
сначала принадлежит одному ядру, а в молекуле втягивается в пространство
между ядрами и связывает их}
\label{fig:H2}
\end{figure}

\begin{glavnoe}
Ковалентная связь~— обобществление \emph{пары} электронов (по одному от
каждого атома, с противоположными спинами), создающих повышенную
электронную плотность между ядрами. Она \vazhno{насыщаема} (число связей
атома ограничено числом неспаренных электронов) и \vazhno{направленна}
(электронные облака перекрываются в определённых направлениях).
\end{glavnoe}

Строгое описание даёт квантовая механика (приближённые решения уравнения
Шрёдингера для молекул), но основные выводы можно получить на языке
квантовых ячеек и форм электронных облаков.

\paragraph{$\sigma$- и $\pi$-связи}
У фтора один неспаренный электрон на $2p$-оболочке, в «гантели»,
вытянутой вдоль одной оси. Два атома F могут образовать связь, только
сблизившись вдоль этой оси (направленность связи!), и тогда облака
перекрываются «лоб в лоб». В молекуле O$_2$, кроме такой пары, есть ещё
пара $p$-электронов с «гантелями», перпендикулярными оси молекулы; они
перекрываются «боком», и плотность концентрируется по обе стороны от оси
(рис.~\ref{fig:sigma-pi}).

\begin{figure}[ht]
\centering
\begin{tikzpicture}[scale=0.8]
  \tikzset{lobeA/.style={fill=fmBlue!30, draw=fmBlue, line width=0.5pt},
           lobeB/.style={fill=fmGreen!30, draw=fmGreen!80!black, line width=0.5pt}}
  % sigma s-s
  \begin{scope}
    \draw[lobeA, fill opacity=0.7] (-0.45,0) circle (0.7);
    \draw[lobeB, fill opacity=0.7] (0.45,0) circle (0.7);
    \fill (-0.45,0) circle (1.5pt); \fill (0.45,0) circle (1.5pt);
    \node[slbl] at (0,-1.3) {$\sigma$ ($s$--$s$)};
  \end{scope}
  % sigma p-p
  \begin{scope}[shift={(3.8,0)}]
    \draw[lobeA, fill opacity=0.7] (-1.0,0) ellipse (0.55 and 0.32) (-0.0,0) ellipse (0.55 and 0.32);
    \draw[lobeB, fill opacity=0.7] (0.1,0) ellipse (0.55 and 0.32) (1.1,0) ellipse (0.55 and 0.32);
    \fill (-0.45,0) circle (1.5pt); \fill (0.55,0) circle (1.5pt);
    \node[slbl] at (0,-1.3) {$\sigma$ ($p$--$p$)};
  \end{scope}
  % pi
  \begin{scope}[shift={(7.6,0)}]
    \draw[lobeA, fill opacity=0.7] (-0.4,0.5) ellipse (0.32 and 0.55) (-0.4,-0.5) ellipse (0.32 and 0.55);
    \draw[lobeB, fill opacity=0.7] (0.4,0.5) ellipse (0.32 and 0.55) (0.4,-0.5) ellipse (0.32 and 0.55);
    \fill (-0.4,0) circle (1.5pt); \fill (0.4,0) circle (1.5pt);
    \draw[fmGray, dashed] (-1.1,0) -- (1.1,0);
    \node[slbl] at (0,-1.3) {$\pi$ ($p$--$p$ боком)};
  \end{scope}
\end{tikzpicture}
\caption{$\sigma$-связь: облака перекрываются вдоль оси молекулы, плотность
максимальна на оси. $\pi$-связь: облака перекрываются боком, плотность
сосредоточена по обе стороны от оси (пунктир)}
\label{fig:sigma-pi}
\end{figure}

\begin{opredelenie}
\term{$\sigma$-связь}~— связь, при которой электронная плотность
сосредоточена вдоль оси, соединяющей ядра; \term{$\pi$-связь}~— по обе
стороны от оси. Сначала всегда образуется $\sigma$-связь, затем, если
хватает валентных электронов, одна или две $\pi$-связи: так получаются
одинарная (2 электрона), двойная (4) и тройная (6 электронов) связи.
\end{opredelenie}

Молекулы из одинаковых атомов (H$_2$, F$_2$, O$_2$) называют
\term{гомоядерными}; связь в них чисто ковалентная.

\paragraph{Молекулы H$_2$O и NH$_3$}
У кислорода две наполовину занятые $2p$-ячейки~— две «гантели» вдоль осей
$y$ и $z$. Атомы водорода присоединяются вдоль этих осей, и можно было бы
ожидать угол между связями $90^\circ$. Но кислород более
электроотрицателен, чем водород, электронная плотность смещена к
кислороду, концы молекулы с водородом заряжены положительно и
отталкиваются; поэтому угол H--O--H увеличен до $104{,}5^\circ$, а
молекула имеет постоянный дипольный момент. Длина связи O--H
$\approx0{,}096$~нм\ispr{В оригинале (с.~337) указано 0,107~нм; длина связи
O--H в молекуле воды составляет 0,096~нм.}~— близко к сумме ковалентных
радиусов кислорода и водорода ($0{,}066+0{,}037=0{,}103$~нм). Расстояние
H--H ($\approx0{,}15$~нм) вдвое больше, чем в молекуле H$_2$: атомы водорода
в молекуле воды непосредственно не связаны. Многоатомную молекулу
полезно рассматривать как совокупность отдельных связей.

Аналогично у азота три наполовину занятые $p$-ячейки, к ним присоединяются
три атома водорода, а взаимное отталкивание водородных концов
увеличивает углы H--N--H с $90$ до $\approx107^\circ$.

\paragraph{Гибридизация}
Цинк химически активен только в возбуждённом состоянии $4s^14p^1$. Хлор
имеет один неспаренный $3p$-электрон. Формально в ZnCl$_2$ одна связь
должна быть типа $3p$--$4s$, другая~— $3p$--$4p$, и они должны различаться.
Опыт показывает, что связи \emph{одинаковы}. Объяснение: в валентном
состоянии электроны находятся не в $s$- и $p$-, а в смешанных
(\term{гибридных}) состояниях, волновые функции которых~— линейные
комбинации $s$- и $p$-функций.

\begin{opredelenie}
\term{Гибридные орбитали}~— волновые функции, составленные из функций
разных типов ($s$, $p$, $d$), нормированные ($\int\psi_i^2\d V=1$),
ортогональные ($\int\psi_i\psi_k\d V=0$) и эквивалентные друг другу.
Они обеспечивают наибольшее перекрывание облаков и самую прочную связь.
\end{opredelenie}

\begin{figure}[ht]
\centering
\begin{tikzpicture}[scale=0.85]
  \tikzset{cen/.style={circle, fill=fmAmber, draw=atomLine, minimum size=5mm, inner sep=0},
           lig/.style={circle, fill=fmGreen!45, draw=atomLine, minimum size=4mm, inner sep=0},
           bond/.style={line width=1.4pt, atomLine!70}}
  \begin{scope}
    \draw[bond] (-1.3,0) -- (1.3,0);
    \node[lig] at (-1.3,0) {}; \node[lig] at (1.3,0) {}; \node[cen] at (0,0) {};
    \node[slbl] at (0,0.55) {$180^\circ$};
    \node[slbl] at (0,-1.5) {$sp$: ZnCl$_2$};
  \end{scope}
  \begin{scope}[shift={(4.2,0)}]
    \foreach \a in {90,210,330}{\draw[bond] (0,0) -- (\a:1.2); \node[lig] at (\a:1.2) {};}
    \node[cen] at (0,0) {};
    \node[slbl] at (0.6,0.35) {$120^\circ$};
    \node[slbl] at (0,-1.5) {$sp^2$: GaCl$_3$};
  \end{scope}
  \begin{scope}[shift={(8.4,0)}]
    \coordinate (c) at (0,0);
    \coordinate (a1) at (0,1.15); \coordinate (a2) at (-1.05,-0.45); \coordinate (a3) at (1.05,-0.45); \coordinate (a4) at (0.25,-0.75);
    \draw[atomLine!35, dashed] (a1)--(a2)--(a3)--cycle (a1)--(a4) (a2)--(a4) (a3)--(a4);
    \foreach \p in {a1,a2,a3}{\draw[bond] (c) -- (\p); \node[lig] at (\p) {};}
    \node[cen] at (c) {};
    \draw[bond] (c) -- (a4); \node[lig] at (a4) {};
    \node[slbl] at (0.85,0.45) {$109{,}5^\circ$};
    \node[slbl] at (0,-1.5) {$sp^3$: CH$_4$, алмаз};
  \end{scope}
\end{tikzpicture}
\caption{Гибридные связи: $sp$~— линейная молекула (угол $180^\circ$),
$sp^2$~— плоская, связи под $120^\circ$, $sp^3$~— тетраэдр, углы
$109{,}5^\circ$}
\label{fig:hybrid}
\end{figure}

\begin{itemize}
\item \emph{$sp$-гибридизация} (Zn в ZnCl$_2$): две эквивалентные связи под
$180^\circ$~— молекула линейна.
\item \emph{$sp^2$} (Ga в GaCl$_3$, конфигурация $4s^14p^2$): три связи в
одной плоскости под $120^\circ$.
\item \emph{$sp^3$} (C в CH$_4$): углерод в основном состоянии
двухвалентен, в возбуждённом~— четырёхвалентен; четыре связи направлены к
вершинам тетраэдра под $109{,}5^\circ$. Связи $sp^3$~— самые прочные из
рассмотренных.
\end{itemize}

\subsubsection{Структуры валентных кристаллов: правило $8-N$}

\begin{keyeq}\label{eq:8-N}
\text{координационное число}=8-N ,
\end{keyeq}
где $N$~— номер группы периодической системы, а 8~— ёмкость $s$- и
$p$-оболочек. $8-N$~— число неспаренных электронов атома, т.\,е. число
соседей, с которыми он должен быть связан, чтобы насытить валентность.
Правило выполняется, если кристалл построен только из нормальных
ковалентных связей (табл.~\ref{tab:8-N}).

\begin{table}[ht]
\centering
\caption{Правило $8-N$ для валентных кристаллов}
\label{tab:8-N}
\small
\begin{tabularx}{\textwidth}{@{}cc>{\raggedright\arraybackslash}p{25mm}>{\raggedright\arraybackslash}X@{}}
\toprule
Группа & к.\,ч. & Элементы & Строение \\
\midrule
VII & 1 & I (и др. галогены) & пары атомов~— молекулы I$_2$ ($d=0{,}266$~нм), связанные силами Ван-дер-Ваальса в цепочки и кристалл \\
VI & 2 & Se, Te; S & Se, Te: бесконечные винтовые цепочки (валентный угол $105^\circ$ и $102^\circ$), гексагональный кристалл; S: замкнутые кольца S$_8$ (S--S $0{,}208$~нм, между молекулами $0{,}36$~нм; 128 атомов в ячейке) \\
V & 3 & As, Sb, Bi & двойные гофрированные слои, каждый атом связан с тремя соседями; слои связаны ван-дер-ваальсовыми силами (ромбоэдрическая решётка) \\
IV & 4 & C (алмаз), Si, Ge, серое Sn & решётка алмаза: каждый атом~— в центре тетраэдра из четырёх соседей \\
\bottomrule
\end{tabularx}
\end{table}

\begin{table}[ht]
\centering
\caption{Межатомные расстояния в двойных слоях мышьяка, сурьмы и висмута:
$b$~— ковалентная связь внутри слоя, $a$~— ближайшее расстояние между
слоями}
\label{tab:AsSbBi}
\small
\begin{tabular}{@{}lccc@{}}
\toprule
Элемент & $a$, нм & $b$, нм & $b/a$ \\
\midrule
As & 0,315 & 0,251 & 0,80 \\
Sb & 0,337 & 0,287 & 0,85 \\
Bi & 0,347 & 0,310 & 0,89 \\
\bottomrule
\end{tabular}
\end{table}
\noindent Чем тяжелее элемент, тем ближе $b/a$ к единице: связь внутри и
между слоями выравнивается, и свойства становятся всё более
металлическими\ispr{В табл.~3.2 оригинала отношение $b/a$ указано как
0,08; 0,085; 0,09~— пропущен порядок: $0{,}251/0{,}315=0{,}80$ и т.\,д.}.

\paragraph{Решётка алмаза}
Её строят так: большой куб~— ГЦК-решётка атомов углерода; куб делят на
8 малых кубиков и в центры четырёх из них (через один) помещают ещё по
атому (рис.~\ref{fig:diamond}). Каждый атом оказывается в центре
тетраэдра из четырёх соседей, связанных с ним ковалентными $sp^3$-связями.
Других связей, кроме ковалентных, в алмазе нет. Такую же структуру имеют
кремний, германий и серое олово.

\begin{figure}[ht]
\centering
\begin{tikzpicture}[cubeproj, scale=3.0, line join=round]
  \CubeEdges[atomLine!50]
  \foreach \p/\q in {
    {(0.25,0.25,0.25)}/{(0,0,0)}, {(0.25,0.25,0.25)}/{(0.5,0.5,0)}, {(0.25,0.25,0.25)}/{(0.5,0,0.5)}, {(0.25,0.25,0.25)}/{(0,0.5,0.5)},
    {(0.75,0.75,0.25)}/{(1,1,0)}, {(0.75,0.75,0.25)}/{(0.5,0.5,0)}, {(0.75,0.75,0.25)}/{(1,0.5,0.5)}, {(0.75,0.75,0.25)}/{(0.5,1,0.5)},
    {(0.75,0.25,0.75)}/{(1,0,1)}, {(0.75,0.25,0.75)}/{(0.5,0,0.5)}, {(0.75,0.25,0.75)}/{(1,0.5,0.5)}, {(0.75,0.25,0.75)}/{(0.5,0.5,1)},
    {(0.25,0.75,0.75)}/{(0,1,1)}, {(0.25,0.75,0.75)}/{(0,0.5,0.5)}, {(0.25,0.75,0.75)}/{(0.5,1,0.5)}, {(0.25,0.75,0.75)}/{(0.5,0.5,1)}}{
    \draw[fmAmber!80!black, line width=1.6pt] \p -- \q;}
  \FCCback[3.6mm] \FCCmid[3.6mm]
  \foreach \p in {(0.25,0.25,0.25),(0.75,0.75,0.25),(0.75,0.25,0.75),(0.25,0.75,0.75)} \node[atomB=3.6mm] at \p {};
  \FCCfront[3.6mm]
\end{tikzpicture}
\caption{Элементарная ячейка алмаза: ГЦК-решётка (голубые атомы) и четыре
атома внутри (оранжевые) в центрах четырёх из восьми малых кубиков; линии~—
ковалентные связи, каждый атом связан с четырьмя соседями по тетраэдру}
\label{fig:diamond}
\end{figure}
